
\subsection{Lemma 2B --- Representation of Functional Dependencies}

\begin{lemma}
\llabel{functionaldependencyrepresentationlemma}
If $\catc$ is a locally finite \newt{RR.5 range} category and $\reqtc$ is a H of instances, if $\catc$ is 
\term{maximally constrained} to the requirement $\reqtc$ then
all functional dependencies $\fundep{H}{f}{g}$  with respect to $\reqtc$ are represented in $\catc$.
\end{lemma}
\begin{proof}

Suppose such a category  $\catc$  that  is 
\term{maximally constrained} to a requirement $\reqtc$ and suppose
\fgsourcediag
in $\catc$ 
and that there is a functional dependency $\fundep{H}{f}{g}$ with respect to $\reqtc$.

From sketch $S$ of \catcw we can construct a sketch $S'$ by formally 
adding a morphism $\qq{h}: b \morph c$
and path equivalences 
\begin{equation}
\label{hquote}
f \circ \qq{h} = g,
\end{equation} 
\begin{equation}
\label{hquotebar}
\overline{\qq{h}}\leq\hat{f}
\end{equation} 
and 
\newt{from which follows, because we are in a RR.5 range category,}
\begin{equation}
\label{hquotehat}
\widehat{\qq{h}}=\hat{g}. 
\end{equation}
\begin{worktt}the proof of this and the proof of the second part
of lemma \lref{fdrangesublemma} are identical. Make a lemma
regarding f,g and h in the earlier basics of range categories section.
The proof is in a hand written note.
\end{worktt}

Let $\catcp$ be the range category generated by $S'$ and
let $I: \catc\morph \catcp$ be the inclusion (range) functor 
generated by the inclusion of $S$ in $S'$. \commentary{I am not immediately sure if the faithfulness of $I$ 
follows from the definition, though it probably does. The faithfulness will surely follow, if we need it, from the existence of the functional dependency and from equational completeness.}
The inclusion functor is a range functor because the restriction and range operators of \catcpw are extensions of those
of \catc.

Now we show that $I$ is consistent with $\reqtc$. 
We need to show that any $D \in \reqtc$
extends unqiuely to $D' :\catcp \morph \FinPar$. Assume such a $D$. 
$D$  extends to $\catcp$ iff there is a  function that we can choose as the value of  for $D'(\qq{h})$  such that $D'(f) \circ D'(\qq{h}) = D'(g)$ i.e such that
$D(f) \circ D'(\qq{h}) = D(g)$ and $\overline{D'(\qq{h})}\leq \widehat{D(f)}$
and so we  extend $D$ to $D'$ 
 by defining $D'(\qq{h})=H_D$. 

\begin{worktt}Why can D be extended to D'?
We have to give a slightly awkward description of D', I think.
It is basically because D' is freely generated. 
\end{worktt} 
This extension to $D'$ is unique because 
lemma \lref{UniquenessOfHD}
$H_D$ is the unique function 
that satisfies $D(f) \circ H_D = D(g)$ and $\overline{H_D}\leq \widehat{D(f)}$.

In the remainder of this proof 
we define a functor $G: \catc \morph \FinPar$ such that
$G(a)$ has a pair of distinct elements, call them $left$ and $right$, say,  that are mapped by $G(f)$ to identical elements
of $G(b)$, i.e. such that $G(f)(left)= G(f)(right)$ and, significantly,
such that $G(g)$ is defined on both $left$ and $right$.   
We  argue that because of  maximal constrainedness that $G$ can be extended to a functor 
$G' : \catcp \morph \FinPar$. 
Because $f \circ \qq{h} = g$ in \catcpw we can argue that $G'(g)$,
and thus $G(g)$,
maps $left$ and $right$ to identical elements of $G(c)$ 
i.e. is such that $G(g)(left)= G(g)(right)$. 
From the particular definition of $G$ this  is then enough 
to ensure that there exists a morphism $h:b \morph c$ in \catcw
such that $f \circ g = h$
 with the required properties to represent the functionality dependency $H$.

From the \newt{RR.5 range} category \catcw we can  construct a range category $\catc_s$ in which all \commentary{we need to talk about the embedding into sticklebacks here\highlight{sticklebacks}} morphisms are split and an inclusion range functor $E: \catcw \morph \catc_s$ which is faithful.   \begin{worktt}Or we could talk about the functor
 $E \circ Hom(a,-)$ elsewhere. BTW: I have sometimes called this $H_a$ but that is confusing. \end{worktt}

Define the functor $F: \catc \morph \FinPar$ be the
composition of the range functor $E$ with the 
coproduct $HomP_{\catc_s}(a,-) + HomP_{\catc_s}(a,-)$
in the functor category $Fin^{\catc}$ and label the injections $L$ and $R$, respectively,
so that for each object $x$ of $\catc$ the diagram
\begin{center}
$
\begin{array}{c p{0.5cm} c p{0.5cm} c  }
\Rnode{h1}{HomP_{\catc_s}(a,x)}  &&\Rnode{Fx}{F(x)}  &&   \Rnode{h2}{HomP_{\catc_s}(a,x)}       
\end{array} 
$
\ncarr{h1}{Fx}
\alabel{L_x}
\ncarr{h2}{Fx}
\blabel{R_x}
\end{center}
is a coproduct in $\FinPar$. 

\begin{newtt}
With $F$ so defined if  $j:x_1 \morph x_2$ in $\catc$
then both squares in this diagram
\begin{equation*}
\begin{array}{c p{1.5cm} c p{1.5cm} c}
\Rnode{TL}{Hom(a,x_1)}  &\ &  \Rnode{TC}{F(x_1)}\ & & \Rnode{TR}{Hom(a,x_1)}\\[1.2cm]
\Rnode{BL}{Hom(a,x_2)}  &\ &  \Rnode{BC}{F(x_2)}\ & & \Rnode{BR}{Hom(a,x_2)}
\end{array}
\begin{arrows}
\ncarr{TL}{TC}\alabel{L_{x_1}}
\ncarr{BL}{BC}\blabel{L_{x_2}}
\ncarr{TR}{TC}\blabel{R_{x_1}}
\ncarr{BR}{BC}\alabel{R_{x_2}}
\ncarr{TL}{BL}\blabel{Hom(a,j)}
\ncarr{TC}{BC}\alabel{F(j)}
\ncarr{TR}{BR}\alabel{Hom(a,j)}
\end{arrows}
\end{equation*}
commute and so if  $k: a \morph x_1$ and if $\overline{k \circ j}=\overline{k}$ then
\begin{equation}
\label{Lfjkequation}
L_{x_2}(k \circ j) = F(j)(L_{x_1}(k))
\end{equation}
and
\begin{equation}
\label{Rfjkequation}
L_{x_2}(k \circ j) = F(j)(L_{x_1}(k))
\end{equation}
\end{newtt}

Now for each object $x$ of $\catc$, we define an equivalence relation $\sim_x$ on $F(x)$ by defining,
for $k_1,k_2:a \morph x$ in $\catc_s$,\commentary{in the first clause of the definition of $\sim_x$ I mean more precisely that $k_1,k_2$ are in the image of the inclusion functor from \catc}\commentary{In fact this is where I am using faithfulness of the mbedding of \catcw in $\catc_s$. Need to draw this out.}

\begin{align*}
L_x(k_1) \sim_x R_x(k_2) & \mbox{ iff $k_1,k_2:a \morph x$ in $\catc$ 
                                     and there exists $k:b \morph x$ in $\catc$ such that $k_1 = f \circ k = k_2$,}\\
L_x(k_1) \sim_x L_x(k_2) & \mbox{ iff $k_1 = k_2$,} \\
R_x(k_1) \sim_x R_x(k_2) & \mbox{ iff $k_1 = k_2$.} \\
\end{align*}
\begin{rationale}
We could  possibly have added to the first of the three clauses above
that $\bar{k}\leq \widehat{f}$ but we don't need to do 
this because at the
end when the rabbit is pulled out of the hat I can take $\widehat{f} \circ k$ as the representation of functional depedency H.
\end{rationale}
\begin{worktt}
Don't lose sight of the fact that $k_1$ and $k_2$ may be sticklebacks.
Does this impact this definition?
\end{worktt}

It is very easy to see that $\sim_x$  is both symmetric and transitive. \commentary{yes very very easy. }

We can show that equivalence is preserved by  $F(j)$ for each morphism $j$, $j: x_1 \morph x_2$ in $\catc$. For assume such a $j$ then 
\commentary{a little bit different to same argument in preparation paper because $F(j)(k1)$ may be undefined.
 I havent described these cases anywhere}
if $y_1,y_2 \in F(x_1)$ such that $y_1 \sim_{x_1} y_2$
then it follows easily by cases and from the definition of $\sim$ that $F(j)(y_1) \sim_{x_2} F(j)(y_2)$.
\begin{worktt}This is all fine but ought to say a little more...
and need to sort out the two earlier sections on quotients.
\end{worktt}

Therefore, using the construction we described earlier, 
we can define 
 $G: \catc \morph \FinPar$  as the quotient functor $\Fquotient$.
By lemma \ref{rangefunctorquotient}, $G$ is a range functor.

In the earlier explanation we said we would use elements $left$ and $right$ of $G(a)$
in the proof; in what follows
the elements
$[L_a(\bar{f}\circ \bar{g})]$ and $[R_a(\bar{f}\circ \bar{g})]$ fulfill the role of $left$ and $right$, respectively.

\begin{newtt}
With $G$ so defined then if $k: a \morph x_1$ and $j:x_1 \morph x_2$ in $\catc$ such that $\overline{k \circ j}=\overline{k}$
then, from the definition of $G(j)$ and using  (\ref{Lfjkequation}), we get
\begin{equation*}G(j)([L_{x_1}(k)])=[F(j)(L_{x_1}(k))] =[L_{x_2}(k \circ j)]
\end{equation*}
 and likewise from the definition of $G(j)$ and using  (\ref{Rfjkequation}), we get
\begin{equation*} 
G(j)([R_{x_1}(k)])=[R_{x_2}(k \circ j)].
\end{equation*} 
From which, since $\overline{g}: a \morph a$ and $f:a \morph b$  
and since $\overline{\overline{g}\circ f} = \overline{g}$ we have 
\begin{equation}
\label{Lfequation}
G(f)([L_a(\bar{g})])=[L_{b}(\bar{g} \circ f)]
\end{equation}
and likewise since $g:a \morph c$   we have 
\begin{equation}
\label{Lgequation}
G(g)([L_a(\bar{g})])=[L_{c}(\bar{g} \circ g)]=[L_{c}(g)]
\end{equation}
\end{newtt}

Now that we have described the functors  $G: \catc \morph \FinPar$ and $I:\catc \morph \catcp$ that is consistent with $\reqtc$
we can use the fact that $\catc$ is maximally constrained to tell us that $G$ extends to a functor 
$G' : \catcp \morph \FinPar$. Since $f \circ \qq{h} = g$ in $\catcp$ then we have
 $G'(f) \circ G'(\qq{h}) = G'(g): G(a) \morph G(c)$ in $\FinPar$.

\begin{worktt}
\begin{align*}
[L_c(g)] 
        &= G(g)([L_a(\bar{g})])       
                && \mbox{by (\ref{Lgequation}),}\\
        &= G'(g)([L_a(\bar{f}\circ \bar{g})])      
                && \mbox{since $G'$ extends $G$,} \\
        &= G'(\qq{h}) (G'(f)([L_a(\bar{g})])) 
                && \mbox{since $G'(f) \circ G'(\qq{h}) = G'(g)$, since $G'$ functorial,}\\
        &= G'(\qq{h}) (G(f)([L_a(\bar{g})])) 
                && \mbox{since $G'$ extends $G$,}  \\
	      &= G'(\qq{h}) ([L_b(\bar{g} \circ f)])       
                &&\mbox{by (\ref{Lfequation}).} 
\end{align*}
\end{worktt}


\begin{worktt} 
Need be careful, above, because $G'(\qq{h})$ is a stickleback.
\end{worktt}
By symmetry we also have
$$[R_c(g)] = G'(\qq{h}) ([R_b(\bar{g} \circ f)])  $$
and because
$$[L_b(\bar{g} \circ f)]=[R_b(\bar{g} \circ f)]$$
we can therefore can conclude that
$$[L_c(g)] = [R_c(g)]$$
and thus that
$$ L_c(g) \sim_c R_c(g)$$

\workt{The rabbit out of the hat:}

It follows from the definition of $\sim_c$ that there exists $k:b \morph c$ in 
$\catc$ such that 
$f \circ k = g$.  
Therefore the functional dependency $\fundep{H}{f}{g}$ 
is represented in $\catc$ 
by the morphism $\widehat{f} \circ k$.
\end{proof}